\section{Results}

Phân tích chính được thực hiện trên các cửa sổ processed PPG hợp lệ dài $60$~s đáp ứng các tiêu chí kiểm soát chất lượng đã định trước. Các kết quả dưới đây lần lượt đánh giá tổ chức động lực học của tín hiệu so với noisy pseudoperiodic null và sau đó so sánh các đặc trưng động lực học giữa trạng thái Awake và Drowsy.

\subsection{Dynamical Organization Beyond a Noisy Pseudoperiodic Null}

Ở cả Awake và Drowsy, PPG gốc nhìn chung khác với PPS theo một pattern nhất quán trên các nhóm đặc trưng động lực học. Về finite-horizon forecastability, PPG gốc có CC cao hơn và NRMSE thấp hơn tập surrogate. DET và $L_{\mathrm{mean}}$ cao hơn, phản ánh cấu trúc tái diễn theo đường chéo mạnh hơn và kéo dài hơn so với PPS; LLE cũng cao hơn, trong khi LAM và TT nhìn chung thấp hơn. Nhìn chung, các khác biệt này cho thấy tín hiệu gốc và noisy pseudoperiodic ensemble không có cùng tổ chức trên prediction, recurrence structure và local trajectory divergence.

Fig.~\ref{fig:rq1_evidence} tổng hợp độ lệch original-minus-PPS theo session cùng tỷ lệ cửa sổ bác bỏ PPS null. Hướng của các khác biệt được quan sát ở cả Awake và Drowsy và nhìn chung được duy trì trên phần lớn các metric, cho thấy pattern không chỉ giới hạn ở một trạng thái hoặc một estimator riêng lẻ. Tỷ lệ bác bỏ ở cấp cửa sổ cho thấy một phần đáng kể các đoạn PPG quan sát được không được noisy pseudoperiodic null đã kiểm tra tái tạo đầy đủ.

Tuy nhiên, tỷ lệ bác bỏ này chỉ mô tả bằng chứng ở cấp cửa sổ đối với tested null và không phải là kiểm định paired giữa Awake và Drowsy. Do đó, RQ1 chỉ hỗ trợ kết luận rằng cả hai trạng thái chứa tổ chức động lực học vượt quá những gì được tái tạo bởi noisy pseudoperiodic model đã kiểm tra. Việc bác bỏ PPS không được diễn giải như bằng chứng trực tiếp cho deterministic chaos và cũng không loại trừ mọi giải thích ngẫu nhiên có thể có.

\begin{figure*}[t]
    \centering

    \begin{minipage}[t]{0.49\textwidth}
        \centering
        \includegraphics[width=\linewidth]{image/statistic_test_rq1/rq1_pps_primary_metrics_60s.pdf}
        \caption*{\textbf{(A)} Session-level deviation between original PPG and median PPS together with window-level PPS rejection fractions across metrics and states.}
    \end{minipage}
    \hfill
    \begin{minipage}[t]{0.49\textwidth}
        \centering
        \includegraphics[width=\linewidth]{image/statistic_test_rq1/rq1_pps_rank_rejection_heatmap_60s.pdf}
        \caption*{\textbf{(B)} Summary of the directional pattern observed for original PPG relative to the PPS null model.}
    \end{minipage}

    \caption{
    Results of pseudoperiodic surrogate testing.
    (A) Original PPG consistently deviated from the PPS ensemble across multiple nonlinear metrics, while a large fraction of windows rejected the noisy pseudoperiodic null in both Awake and Drowsy states.
    (B) The resulting directional pattern indicates higher forecastability, stronger diagonal recurrence organization, and greater local trajectory divergence in the original PPG, together with lower laminarity-related measures relative to PPS.
    }
    \label{fig:rq1_evidence}
\end{figure*}

\subsection{Awake--Drowsy Differences in PPG Dynamics}

Table~\ref{tab:awake_drowsy_statistics} summarizes the paired effects of the seven metrics in the primary $60$~s analysis, while Fig.~\ref{fig:awake_drowsy_paired} illustrates the session-level changes from Awake to Drowsy. A coordinated pattern was observed across the headline metrics, including lower Mean CC, higher Mean NRMSE, lower DET, and lower LLE in Drowsy. For finite-horizon forecastability, Mean CC decreased with a median paired difference of $\Delta=-0.0339$ and $95\%$ CI $[-0.0595,-0.0043]$, whereas Mean NRMSE increased with $\Delta=+0.0294$ and $95\%$ CI $[0.0047,0.0530]$. Both remained supported after BH-FDR correction, with $q_{\mathrm{BH}}=0.0376$ for Mean CC and $q_{\mathrm{BH}}=0.0255$ for Mean NRMSE, indicating lower finite-horizon forecastability in Drowsy.

\begin{figure*}[t]
    \centering
    \includegraphics[width=\textwidth]{image/paired_comparison_rq2/paired_awake_drowsy_7metrics_final.pdf}
    \caption{
    Paired Awake--Drowsy comparison of the seven nonlinear dynamical metrics obtained from the primary $60$-s processed-PPG analysis. Each gray line connects the Awake and Drowsy values from the same session. Blue and orange markers denote Awake and Drowsy states, respectively. The panels show (a) CC, (b) NRMSE, (c) DET, (d) $L_{\mathrm{mean}}$, (e) LAM, (f) TT, and (g) LLE.
    }
    \label{fig:awake_drowsy_paired}
\end{figure*}

Về recurrence organization, DET giảm với $\Delta=-0.0186$, $95\%$ CI $[-0.0364,-0.0069]$ và $q_{\mathrm{BH}}=0.0376$, cho thấy tổ chức tái diễn theo đường chéo thấp hơn ở Drowsy tại cấu hình tái dựng nominal. $L_{\mathrm{mean}}$ cũng giảm theo cùng hướng với $\Delta=-0.0630$ và CI không chứa $0$, nhưng không còn được hỗ trợ sau BH-FDR ($q_{\mathrm{BH}}=0.0816$). Ngược lại, LAM và TT có xu hướng tăng, song confidence intervals đều chứa $0$ và các kiểm định tương ứng không được hỗ trợ sau hiệu chỉnh multiple comparisons. Vì vậy, các thay đổi liên quan đến vertical hoặc laminar recurrence được xem là secondary và kém nhất quán hơn so với thay đổi của DET.

Về local trajectory divergence, LLE giảm với $\Delta=-0.0376~\mathrm{s}^{-1}$, $95\%$ CI $[-0.0537,-0.0230]~\mathrm{s}^{-1}$, $r_{\mathrm{rb}}=-0.810$ và $q_{\mathrm{BH}}=0.0050$, cho thấy mức phân kỳ quỹ đạo cục bộ ước lượng thấp hơn ở Drowsy. Nhìn chung, kết quả RQ2 cho thấy trạng thái Drowsy liên quan đến sự tái tổ chức phối hợp của PPG trên ba thuộc tính bổ sung: finite-horizon forecastability giảm, tổ chức tái diễn theo đường chéo giảm và mức phân kỳ quỹ đạo cục bộ ước lượng giảm. Các metric còn lại cung cấp bối cảnh bổ sung nhưng không cho thấy mức hỗ trợ thống kê nhất quán tương đương với bốn headline findings.

\begin{table}[t]
    \centering
    \caption{Paired Awake--Drowsy statistical comparison for the primary $60$-s analysis. The paired difference is defined as $\Delta=\mathrm{Drowsy}-\mathrm{Awake}$.}
    \label{tab:awake_drowsy_statistics}
    \footnotesize
    \setlength{\tabcolsep}{3.2pt}
    \renewcommand{\arraystretch}{1.08}

    \begin{tabular}{lcccc}
        \toprule
        \textbf{Metric}
        & \textbf{Median $\Delta$}
        & \textbf{95\% CI}
        & $\mathbf{r_{rb}}$
        & $\mathbf{q_{BH}}$ \\
        \midrule

        CC
        & $-0.0339$
        & $[-0.0595,-0.0043]$
        & $-0.581$
        & $\mathbf{0.0376}$ \\

        NRMSE
        & $+0.0294$
        & $[0.0047,0.0530]$
        & $+0.667$
        & $\mathbf{0.0255}$ \\

        DET
        & $-0.0186$
        & $[-0.0364,-0.0069]$
        & $-0.581$
        & $\mathbf{0.0376}$ \\

        $L_{\mathrm{mean}}$
        & $-0.0630$
        & $[-0.1247,-0.0139]$
        & $-0.486$
        & $0.0816$ \\

        LAM
        & $+0.0165$
        & $[-0.0044,0.0635]$
        & $+0.381$
        & $0.1667$ \\

        TT
        & $+0.0095$
        & $[-0.0034,0.0160]$
        & $+0.333$
        & $0.2024$ \\

        LLE
        & $-0.0376$
        & $[-0.0537,-0.0230]$
        & $-0.810$
        & $\mathbf{0.0050}$ \\

        \bottomrule
    \end{tabular}

    \vspace{1mm}

    \begin{minipage}{0.97\columnwidth}
        \footnotesize
        Bold values indicate metrics supported after Benjamini--Hochberg false-discovery-rate correction ($q_{\mathrm{BH}}<0.05$).
    \end{minipage}
\end{table}

\subsection{Robustness and Inferential Scope}

Mean CC, DET và LLE giữ hướng giảm, còn Mean NRMSE giữ hướng tăng ở Drowsy tại cả bốn độ dài cửa sổ $30$, $60$, $120$ và $180$~s (Fig.~\ref{fig:window_robustness}). Các hướng hiệu ứng được duy trì nhất quán hơn so với độ lớn hiệu ứng. Các ước lượng ở $30$~s có độ bất định lớn hơn. Từ $60$ đến $180$~s, các hiệu ứng nhìn chung ổn định hơn, mặc dù DET tại $180$~s vẫn cho khoảng bất định tương đối rộng. Trong dataset này, $60$~s là một lựa chọn cân bằng thực tế giữa thời lượng quan sát ngắn và độ ổn định của ước lượng, không phải độ dài tối ưu phổ quát. Cả bốn hướng cũng được giữ dưới các label rules P0, T30, T60, S3 và S5 (Table~\ref{tab:label_sensitivity}), cho thấy pattern quan sát được không bị chi phối chủ yếu bởi các cửa sổ gần chuyển tiếp hoặc các episode trạng thái ngắn trong phạm vi đã khảo sát. Trên bốn label rules thay thế, cả $16/16$ headline effects đều giữ hướng kỳ vọng và đạt $q_{\mathrm{BH}}<0.05$, dù bootstrap CI chứa $0$ ở T60--Mean CC và T30--LLE; S5 còn $19$ paired sessions so với $20$ ở các rule còn lại.

\begin{figure*}[t]
    \centering
    \includegraphics[width=0.78\textwidth]{image/robustness_window/robust_window_primary_metrics.pdf}
    \caption{
    Robustness of the primary Awake--Drowsy effects across window lengths. Paired effects ($\Delta=\mathrm{Drowsy}-\mathrm{Awake}$) are shown for (a) CC, (b) NRMSE, (c) DET, and (d) LLE using $30$, $60$, $120$, and $180$ s windows. Error bars represent the corresponding $95\%$ bootstrap confidence intervals. All four headline metrics preserve the same direction of change across window lengths, whereas the $30$-s estimates exhibit greater uncertainty.
    }
    \label{fig:window_robustness}
\end{figure*}

\begin{table}[t]
\centering
\caption{Sensitivity of the primary Awake--Drowsy effects to alternative label definitions. Values are median paired differences $\Delta=\mathrm{Drowsy}-\mathrm{Awake}$. P0 denotes the primary labeling; T30/T60 exclude samples within $\pm30/\pm60$ s of state transitions; S3/S5 retain only continuous state episodes lasting at least $3/5$ min.}
\label{tab:label_sensitivity}
\footnotesize
\setlength{\tabcolsep}{4.5pt}
\renewcommand{\arraystretch}{1.08}
\begin{tabular}{lccccc}
\toprule
\textbf{Metric} & \textbf{P0} & \textbf{T30} & \textbf{T60} & \textbf{S3} & \textbf{S5} \\
\midrule
Mean CC
& -0.0339 & -0.0482 & -0.0251 & -0.0361 & -0.0337 \\
Mean NRMSE
& +0.0294 & +0.0406 & +0.0278 & +0.0357 & +0.0367 \\
DET
& -0.0186 & -0.0183 & -0.0167 & -0.0189 & -0.0213 \\
LLE
& -0.0376 & -0.0322 & -0.0393 & -0.0344 & -0.0321 \\
\bottomrule
\end{tabular}
\end{table}

Các metric có độ bền khác nhau khi thay đổi tham số (Table~\ref{tab:parameter_sensitivity}). Mean CC, Mean NRMSE và LLE giữ hướng trong grid delay và dimension đã kiểm tra. DET giữ hướng âm tại $m=7$, $8$ và $9$, nhưng nhạy rõ với embedding delay: hiệu ứng âm tại $\tau=0.16$~s tiến gần $0$ và đổi sang dương tại $\tau=0.12$ và $0.20$~s, nên kết quả DET cần được diễn giải gắn với cấu hình nominal. Đối với unknown repeated-session dependence, phân tích bổ sung sau Version~1 lấy mẫu $100{,}000$ cách ghép $20$ sessions thành $10$ hypothetical two-session clusters, do liên kết participant--session thực tế của $10$ người không còn khả dụng. Phân tích này sử dụng một baseline Simplex riêng đã được freeze cho sensitivity analysis, khác kết quả primary trong Table~\ref{tab:awake_drowsy_statistics}. Median hiệu ứng của cả bốn headline metrics giữ hướng trong toàn bộ các cách ghép giả định đã lấy mẫu, nhưng độ lớn và mức hỗ trợ suy luận thay đổi. Tỷ lệ đạt $q_{\mathrm{BH}}<0.05$ là $19.36\%$, $27.82\%$, $22.59\%$ và $79.69\%$ cho Mean CC, Mean NRMSE, DET và LLE, tương ứng. Phân tích này đánh giá độ nhạy với cấu trúc cluster chưa biết, không thay thế suy luận primary ở cấp session. Các cluster giả định không khôi phục danh tính participant, chứng minh tính độc lập giữa các sessions hay xác lập suy luận cấp participant.

\begin{table}[t]
\centering
\caption{Sensitivity of the primary effects to phase-space reconstruction parameters. Values are median paired differences $\Delta=\mathrm{Drowsy}-\mathrm{Awake}$. The nominal Mean CC and Mean NRMSE values are the reconstruction-sensitivity estimates reported in Version 1 and differ from the primary estimates in Table~\ref{tab:awake_drowsy_statistics}; they are retained for within-grid sensitivity comparisons and do not replace the primary estimates.}
\label{tab:parameter_sensitivity}
\footnotesize
\setlength{\tabcolsep}{3.5pt}
\renewcommand{\arraystretch}{1.08}
\begin{tabular}{lccc|ccc}
\toprule
& \multicolumn{3}{c}{\textbf{Delay $\tau$ (s), $m=8$}}
& \multicolumn{3}{c}{\textbf{Dimension $m$, $\tau=0.16$ s}} \\
\textbf{Metric}
& 0.12 & 0.16 & 0.20
& 7 & 8 & 9 \\
\midrule
Mean CC
& -0.0317 & -0.0277 & -0.0316
& -0.0307 & -0.0277 & -0.0251 \\
Mean NRMSE
& +0.0321 & +0.0319 & +0.0296
& +0.0377 & +0.0319 & +0.0315 \\
DET
& +0.0025 & -0.0186 & +0.0015
& -0.0145 & -0.0186 & -0.0200 \\
LLE
& -0.0361 & -0.0376 & -0.0464
& -0.0285 & -0.0376 & -0.0247 \\
\bottomrule
\end{tabular}
\end{table}

Ở lớp hỗ trợ exploratory, symbolic dynamics cho thấy $0V\uparrow$ và $2V\downarrow$ tại $60$, $120$ và $180$~s, trong khi 1V gần $0$ và kém nhất quán hơn (Table~\ref{tab:symbolic_results}; chênh lệch tính bằng điểm phần trăm). Pattern này cung cấp một lớp bằng chứng bổ sung về sự thay đổi động lực học theo trạng thái, nhưng không có hỗ trợ thống kê nhất quán qua các cấu hình và không được diễn giải như phép đo trực tiếp hoạt tính sympathetic hoặc parasympathetic. Nhìn chung, trong các kiểm tra đã thực hiện, hướng của các primary effects ổn định hơn độ lớn và mức hỗ trợ thống kê của chúng. Mức độ robustness phụ thuộc từng metric: Mean CC, Mean NRMSE và LLE duy trì hướng tốt hơn trong parameter grid đã khảo sát, trong khi DET nhạy rõ hơn với embedding delay. Các kết quả repeated-session sensitivity tiếp tục hỗ trợ sự ổn định về hướng, nhưng không loại bỏ bất định liên quan đến participant--session dependence.

\begin{table}[t]
    \centering
    \caption{Awake--Drowsy symbolic dynamics across window lengths. Paired differences are defined as $\Delta=\mathrm{Drowsy}-\mathrm{Awake}$.}
    \label{tab:symbolic_results}
    \footnotesize
    \setlength{\tabcolsep}{4pt}
    \renewcommand{\arraystretch}{1.10}

    \begin{tabular}{lcccc}
        \toprule
        \textbf{Metric}
        & \textbf{60 s}
        & \textbf{120 s}
        & \textbf{180 s}
        & \textbf{Same direction} \\
        \midrule

        0V
        & $+4.04$
        & $+2.32$
        & $+4.58$
        & $17/20$ \\

        1V
        & $+0.41$
        & $+0.03$
        & $-0.06$
        & $11/20$ \\

        2V
        & $-0.91$
        & $-2.30$
        & $-2.67$
        & $14/20$ \\

        \bottomrule
    \end{tabular}
\end{table}
